\documentclass[11pt,a4paper]{article}

\usepackage{polyglossia}
\setmainlanguage{brazil}

\usepackage{fsorg-manual}

\graphicspath{{figures/pt_BR/}}

\hypersetup{
  pdftitle={FS Organizer · Manual do usuário},
  pdflang={pt-BR}}

\fsorgrunningheadis{FS Organizer · Manual do usuário}

\begin{document}

\frenchspacing

\fsorgtitle
  {Manual do usuário}
  {Versão \input{../VERSION.txt}}
  {Distribuído sob a GNU General Public License versão 2.}
  {A tipografia é a Archivo, de Omnibus-Type, sob a SIL Open Font License 1.1.}

\thispagestyle{empty}
\tableofcontents

\section{O que este programa faz}

O FS Organizer ativa e desativa addons do Microsoft Flight Simulator sem mover
os arquivos deles.

Você guarda os seus addons numa pasta fora do simulador, que o programa chama de
\ui{biblioteca}. As subpastas dela são as \ui{categorias}. Ativar um addon cria
um link para ele numa pasta que o simulador lê, chamada \ui{destino}. Desativar
remove o link. Os arquivos ficam na biblioteca o tempo todo, então ativar ou
desativar um addon, mesmo grande, é instantâneo.

\subsection{O que o programa não guarda}

O programa nunca anota quais addons estão ativados. A cada varredura, ele olha
os destinos no disco e mostra o que está lá. Se você mexer nas pastas à mão, ou
se o instalador de outro programa mexer, clique em \ui{Atualizar} e a lista
volta a bater com o que o simulador vê.

\subsection{Perfil}

Um \ui{perfil} é uma instalação do simulador, com os destinos e as bibliotecas
dela. Quem tem o Flight Simulator 2020 e o 2024 no mesmo computador tem dois
perfis. O programa só altera o perfil em uso, que você escolhe em \ui{Opções}.

\subsection{Tipo de link}

Em \ui{Opções}, na aba \ui{Links}, você escolhe como os links são criados.

\begin{description}
  \item[Junção de diretório] Funciona entre discos locais sem precisar de
    administrador. É o padrão, e é a que você deve usar a menos que a biblioteca
    esteja numa pasta de rede.
  \item[Link simbólico] Só é necessário para biblioteca numa pasta de rede. O
    Windows só permite com o Modo de Desenvolvedor ligado ou com direitos de
    administrador, e o programa avisa quando nenhum dos dois está disponível.
\end{description}

Trocar o tipo vale só para os links novos. Os que já existem continuam como
estão.

\section{Presets}

Um preset é um conjunto de addons ativados, guardado com um nome. O mesmo preset
pode ser aplicado de três jeitos, e cada um dá um resultado diferente, então
leia o plano antes de aplicar.

\subsection{O que um preset guarda}

Para criar um preset, ative os addons que você quer e clique em \ui{Novo a
partir dos ativados}.

A coluna \ui{Conteúdo} mostra quantos addons o preset cita e em quantas
categorias eles estão. A coluna \ui{Atualizado} mostra quando ele foi salvo pela
última vez.

Um preset também pode controlar as entradas de inicialização do simulador
(seção~\ref{sec:presets-inicializacao}).

\subsection{Preset igual ao atual}

A coluna \ui{Ao aplicar} mostra quantos addons mudariam se você aplicasse o
preset agora, no modo escolhido em \ui{Aplicar como}.

\ui{Igual ao atual} aparece ao lado de um preset quando os addons ativados são
exatamente os que ele cita. Essa comparação é sempre feita como Substituir: se
houver mais um addon ativado que o preset não cita, ele deixa de ser igual,
porque Substituir desativaria esse addon. Um preset que você acabou de salvar já
aparece como igual ao atual, mesmo sem nunca ter sido aplicado.

Como a contagem segue o modo escolhido e a marca não, um preset pode estar igual
ao atual e ainda mostrar mudanças. No modo Desativar, por exemplo, ele
desativaria todos os addons que cita.

\subsection{Os três modos de aplicação}

\begin{description}
  \item[Substituir] Ativa os addons do preset e desativa todo o resto.
  \item[Acumular] Ativa os addons do preset e deixa o resto como está.
  \item[Desativar] Desativa os addons do preset. Neste modo, as entradas que o
    preset marca como desativadas são ignoradas.
\end{description}

Antes de aplicar, o painel \ui{Plano} divide o resultado em cinco grupos:

\begin{itemize}
  \item \ui{A ativar}: addons do preset que estão desativados agora.
  \item \ui{A desativar}: addons que o preset desativa, ou que o Substituir
    desativa porque o preset não os cita.
  \item \ui{Já estão como o preset pede}: addons que ficam como estão.
  \item \ui{Não encontrados na biblioteca}: addons que o preset cita e que não
    estão mais na biblioteca. O resto do preset é aplicado mesmo assim.
  \item \ui{Pedidos, mas não aplicados}: entradas de inicialização que o preset
    pede e que não puderam ser alteradas.
\end{itemize}

\shot{03b-presets-plan.png}{O plano do preset ``Voo curto'' no modo Substituir.
A frase abaixo dos grupos diz que os 4 addons fora do preset fazem parte dos 4
que serão desativados.}

\subsection{O que Substituir desativa por omissão}

No modo Substituir, os addons que estão ativados agora e que o preset não cita
são desativados. Eles aparecem em \ui{Desativados pelo Substituir}, e \ui{Ver
quais} abre a lista.

Eles já entram na conta do plano. Se o plano diz que desativa 12 e a lista tem 5
nomes, são 12 addons desativados no total: esses 5, mais 7 que o próprio preset
desativa.

\subsection{O preset de retorno}

Logo antes de aplicar um preset, o programa guarda os addons ativados naquele
momento como o preset de retorno do perfil. \ui{Voltar ao conjunto anterior}
aplica esse preset. Se ele não puder ser guardado, por exemplo porque a pasta de
presets está protegida contra escrita, nada é aplicado.

Isso é diferente do \ui{Desfazer última alteração}:

\begin{description}
  \item[Desfazer última alteração] Desfaz o que você acabou de aplicar. Deixa de
    estar disponível quando você importa um addon, mexe na quarentena ou altera
    uma categoria.
  \item[Voltar ao conjunto anterior] Aplica o preset de retorno como qualquer
    outro preset. Continua disponível depois de importações e de operações na
    quarentena.
\end{description}

\subsection{Presets e entradas de inicialização}
\label{sec:presets-inicializacao}

Na aba \ui{Inicialização} do preset, uma caixa faz o preset controlar também os
programas que o simulador abre junto consigo. Marcar a caixa guarda as entradas
de inicialização ativadas naquele momento. Depois dá para ativar ou desativar
cada uma ali.

Se o gerenciamento de inicialização estiver desativado em \ui{Opções}, as
entradas do preset não são aplicadas, e o programa diz quantas ficaram de fora.
Os addons são aplicados normalmente.

Se uma entrada que o preset cita não estiver mais no arquivo do simulador, o
programa a lista. Ele nunca acrescenta entradas de inicialização.

\section{A primeira configuração}

Na primeira vez que você abre o programa, ele procura as instalações do
simulador e pede duas coisas:

\begin{description}
  \item[O simulador] Escolha na lista a instalação deste perfil, ou escolha uma
    pasta você mesmo. Se a pasta não parecer a pasta \where|Community| do
    simulador, o programa avisa e usa assim mesmo.
  \item[A biblioteca] Escolha a pasta fora do simulador onde ficam os seus
    addons. As subpastas dela viram categorias.
\end{description}

Ao adicionar uma biblioteca com caminho longo, o programa mostra quantos
caracteres o caminho já ocupa e compara com um caminho curto. O Windows recusa
algumas operações, a Lixeira entre elas, em caminhos acima de 260 caracteres. O
programa não bloqueia nada: ele avisa quando um caminho atrapalha, e os
\ui{Diagnósticos} mostram o caminho mais longo de cada biblioteca.

Se você usava o MSFS Addons Linker, abra as \ui{Opções} e use \ui{Importar do
MSFS Addons Linker}. Ele adiciona as bibliotecas, as categorias e os presets que
ainda não existem aqui. Nenhum arquivo é movido ou excluído.

\section{Biblioteca}

A aba \ui{Biblioteca} mostra os seus addons numa árvore, agrupados por
categoria. A caixa de cada linha ativa ou desativa o addon.

\shot{01-library.png}{A biblioteca, com as categorias abertas e o painel do
addon selecionado à direita.}

\subsection{As colunas}

\begin{description}
  \item[Addon] O nome da pasta.
  \item[Versão] A versão declarada no \where|manifest.json| do addon.
  \item[Ativado] Se o addon tem um link num destino agora.
  \item[Destino] Onde o link está, e se o destino é fixado.
  \item[Link] Se o link funciona. Quando a pasta dele não existe, o link está
    quebrado.
\end{description}

Uma linha pode ter uma de três marcas: \ui{Link quebrado}, \ui{Em conflito} ou
\ui{Duas cópias}. As seções sobre destinos e quarentena explicam cada uma.

\subsection{Ativar e desativar em massa}

Selecione várias linhas e use \ui{Ativar selecionados} ou \ui{Desativar
selecionados}. Com mais de 10 addons, o programa pergunta antes. O
\ui{Desfazer última alteração} desfaz o lote inteiro.

Se outro addon seu já usa o mesmo nome de pasta no destino, o programa oferece
trocar os dois: desativa aquele e ativa o seu num passo só, que se desfaz de
uma vez.

\subsection{Categorias}

Clique com o botão direito na árvore para criar, renomear ou excluir uma
categoria, ou para mover addons entre categorias. Só categoria vazia pode ser
excluída.

\ui{Sugerir categorias} olha os nomes das pastas e os \where|manifest.json| e
propõe para onde mover cada addon. As regras confiáveis já vêm marcadas. A regra
de livery erra com frequência, então vem desmarcada: revise antes de aplicar.

O programa põe um arquivo oculto \where|.fsorg-category| em cada pasta de
categoria que conhece, para que uma categoria vazia continue reconhecida. Ele
marca as pastas quando você adiciona a biblioteca, e \ui{Categorias}, nas
opções, marca de novo ou remove todos os marcadores.

Uma pasta sem addons dentro e sem \where|manifest.json| próprio conta como
addon, então dá para ativar, mover e excluir como qualquer outro. É assim que
aparecem os pacotes convertidos que o simulador grava e as bibliotecas de
objetos de cenário que vêm sem manifesto.

\subsection{Fixar o destino}

O link de um addon fica no destino padrão do perfil. \ui{Usar sempre}, seguido
de um destino, fixa um addon ou uma categoria inteira nesse destino. Se um destino
fixado for removido do perfil, o programa avisa e usa o destino padrão até você
escolher outro. A fixação em si continua guardada.

\subsection{Reapontar}

Se um link aponta para o lugar errado, ou para uma pasta que não existe mais,
\ui{Reapontar para a biblioteca} corrige. Só o link muda; os arquivos do addon
não são tocados.

\section{Destinos}

A aba \ui{Destinos} lista tudo o que está nas pastas que o simulador lê,
inclusive pastas que você não pôs lá pelo FS Organizer. Cada entrada recebe uma
classificação:

\shot{02-community.png}{Os destinos, com o filtro de classificação no topo e a
contagem de cada tipo embaixo.}

\begin{description}
  \item[Gerenciada] O simulador carrega da sua biblioteca. É o estado normal.
  \item[Externa] Pasta comum instalada por outro programa, ainda fora de uma
    biblioteca.
  \item[Divergente] O programa que instalou o addon recolocou a própria cópia,
    então existem duas cópias.
  \item[Sumida] A cópia da biblioteca sumiu, removida pelo outro programa. Não
    há o que reparar: os arquivos não estão mais neste computador.
  \item[Quebrada] O link aponta para uma pasta que não existe.
  \item[Indisponível] O disco não está conectado. Ele pode voltar, então o
    programa não oferece limpeza.
  \item[Não gerenciada] Pasta comum, ainda fora de uma biblioteca.
  \item[Duplicada] Com link em mais de um destino.
  \item[Substituída] Algo gravou uma pasta comum por cima de um link que este
    programa criou. O simulador carrega essa pasta, e ativar, presets e a busca
    pelo culpado não alcançam mais a cópia da biblioteca.
\end{description}

O filtro no topo mostra uma classificação por vez, e \ui{Selecionar tudo o que
aparece} seleciona exatamente o que está na tela.

\subsection{Importar a pasta de outro programa}

\ui{Importar esta pasta} move a pasta para a sua biblioteca e deixa um link no
lugar dela, e o addon passa a ser seu.

O outro programa não vai saber do link. A próxima atualização dele pode gravar
dentro do link, ou trocá-lo por uma pasta comum e deixar você com duas cópias.
Quando isso acontece, a entrada aparece como \ui{Divergente} e o programa
pergunta qual cópia manter.

Para devolver o addon ao programa que o instalou, use \ui{Excluir da
biblioteca}.

\section{Importar}

Importar copia uma pasta comum para a biblioteca e deixa um link onde ela
estava. A original só é removida depois que a cópia é conferida. O histórico
registra cada um dos cinco passos:

\shot{27-community-import.png}{O diálogo de importação: quanto vai ser copiado,
para qual biblioteca e categoria, e onde o addon vai ficar.}

\begin{enumerate}
  \item Copiar para uma pasta temporária.
  \item Conferir a cópia.
  \item Pôr a cópia no lugar.
  \item Remover a pasta original.
  \item Criar o link no destino.
\end{enumerate}

Enquanto a cópia corre, uma janela mostra a pasta atual, o progresso e o botão
\ui{Cancelar}. Cancelar não exclui nada: o que já foi copiado fica guardado, e a
importação pode ser retomada depois.

\subsection{Verificação depois de copiar}

Na aba \ui{Links} das \ui{Opções}, você escolhe como a cópia é conferida.

\begin{description}
  \item[Por estrutura] Compara a quantidade e o tamanho dos arquivos. É o padrão,
    e é rápida.
  \item[Por hash] Compara o conteúdo de cada arquivo. Pega mais diferenças, mas
    deixa a importação várias vezes mais lenta, principalmente com muitos
    arquivos pequenos.
\end{description}

Se a cópia não conferir, a original não é removida.

\subsection{Importações incompletas}

Se uma importação for interrompida, na próxima vez que o programa abrir ele
oferece três opções: retomar a importação, descartar a cópia parcial ou deixar
como está. A pasta original não é tocada em nenhum dos casos.

Se a operação interrompida era uma resolução de conflito, a cópia parcial só
pode ser descartada, porque as duas cópias continuam no lugar.

\section{Quarentena}

Quando você resolve um conflito entre duas cópias de um addon, a cópia que você
não mantém vai para a quarentena em vez de ser excluída. Os itens ficam lá até
você restaurar ou descartar.

Para cada item, o programa grava de onde ele veio, num arquivo ao lado do item.
Esse registro não depende do histórico, então dá para restaurar um item mesmo
que o histórico se perca. A tabela mostra as duas fontes, \ui{O registro diz} e
\ui{O Histórico diz}, e avisa quando elas discordam.

\shot{05-quarantine.png}{A quarentena, com a origem de cada item e o aviso onde o
registro e o histórico discordam.}

\ui{Restaurar} devolve cada pasta ao lugar de onde veio. Se já houver algo lá, o
diálogo mostra as duas versões e você decide; substituir move a pasta que está
lá para a quarentena. \ui{Descartar} exclui o item definitivamente, depois de
perguntar.

\section{Excluir addon da biblioteca}

\ui{Excluir} oferece três jeitos de tirar um addon da biblioteca, e avisa antes
quando um deles não vai funcionar.

\begin{description}
  \item[Mover para a Lixeira] Dá para recuperar o addon pela Lixeira. Essa opção
    é recusada quando o addon não cabe na Lixeira daquele disco, ou quando tem
    um caminho maior que os 260 caracteres que a Lixeira aceita. O Windows pode
    descartar itens antigos da Lixeira para abrir espaço.
  \item[Excluir definitivamente] Não pode ser desfeito.
  \item[Devolver ao programa que o instalou] Só aparece para addon importado da
    pasta de outro programa. Os arquivos voltam para aquela pasta e saem da
    biblioteca.
\end{description}

Se o addon estiver ativado, ele é desativado antes.

\section{Simulador}

A aba \ui{Simulador} tem dois painéis, e os dois alteram arquivos do próprio
simulador. Antes de alterar um deles, o programa guarda uma cópia de segurança
ao lado. Nenhum dos painéis altera nada com o simulador aberto. Os dois começam sem
gerenciamento, e enquanto for assim o programa não lê nem grava o arquivo:
clique em \ui{Gerenciar as entradas de inicialização} ou em \ui{Gerenciar a
lista de pacotes} para ligar cada um.

\subsection{Entradas de inicialização}

O arquivo \where|EXE.xml| do simulador, ao lado do \where|UserCfg.opt|, lista os
programas que ele abre junto consigo. O painel \ui{Entradas de inicialização}
mostra esses programas e deixa você ativar ou desativar cada um sem editar XML.
O programa só ativa ou desativa entradas que já existem. Ele nunca acrescenta,
remove nem reordena.

\shot{19-simulator-startup.png}{As entradas de inicialização, com o programa que
cada uma abre e o interruptor dela.}

Quando você desativa um addon que tem uma entrada de inicialização, o programa
avisa: sem isso, o simulador seguiria tentando abrir um programa que não está
mais lá. Dá para desativar os dois num passo só, que se desfaz de uma vez.

\subsection{Lista de pacotes}

O simulador traz aeroportos próprios. Quando um aeroporto seu cobre o mesmo
lugar que um deles, o painel \ui{Lista de pacotes} mostra o par e deixa você
desativar o aeroporto do simulador.

Isso altera um valor na lista de pacotes do simulador. Nada é acrescentado,
removido ou reordenado. Uma entrada que não tem esse valor fica como está,
porque um arquivo inválido faz o simulador zerar a lista inteira.

Quando dois addons seus cobrem o mesmo aeroporto, o programa só mostra o par.
Para ficar com um só, desative o outro na aba \ui{Biblioteca}.

\section{Documentos}

A aba \ui{Documentos} encontra os PDFs que vêm com os seus addons e os abre
dentro do programa. \ui{Analisar a biblioteca} procura PDFs em cada addon. O
resultado fica salvo, então as próximas análises são rápidas.

Os PDFs são separados em \ui{Documentos} e \ui{Cartas}. As cartas são agrupadas
por aeroporto, e também por tipo quando o addon traz um catálogo delas. Um addon
pode trazer um manual e cento e cinquenta cartas, por isso cada grupo mostra a
contagem primeiro.

O leitor tem índice, busca, ajuste à largura e marcadores. Um marcador sem nome
recebe o número da página. \ui{Abrir numa janela à parte} tira o leitor da aba.

\shot{26-documents.png}{A aba de documentos: as cartas agrupadas por aeroporto à
esquerda, o leitor no meio e o índice do documento à direita.}

Na carta, a roda do mouse dá zoom e arrastar move a página. No documento, os
dois começam desativados. Os dois botões da barra ativam e desativam cada um.

\section{Diagnósticos}

A aba \ui{Diagnósticos} mostra o que o programa encontrou no disco. Cada seção
diz quando foi medida, e dá para medir de novo a qualquer momento.

\shot{09-diagnostics-size.png}{A seção de tamanho em disco, com o total por
categoria e por addon, e a hora da medição.}

\begin{description}
  \item[Entradas dos destinos] Quantas entradas existem de cada tipo, e o
    caminho mais longo de cada biblioteca.
  \item[Quebradas ou indisponíveis] Links quebrados e entradas num disco
    desconectado, com um botão para reparar os links quebrados.
  \item[Quarentena] Quanto espaço a quarentena ocupa.
  \item[Tamanho em disco] Quanto cada categoria e cada addon ocupam. Pode levar
    alguns segundos; \ui{Parar} para depois do addon atual.
  \item[Aeroportos no cenário] Analisa o cenário de cada addon e agrupa os
    addons pelo que encontra: código de aeroporto, registro de aeroporto
    ilegível, dados de navegação em vez de cenário, ou nenhum registro de
    aeroporto.
  \item[Módulos carregados] Lê o relatório que o simulador grava depois de um
    carregamento demorado. O simulador não informa o tempo de carregamento por
    pacote, então aqui aparece o módulo que cada pacote carregou e quanta
    memória ele ocupa.
\end{description}

\subsection{Encontrar o culpado}

Quando o simulador trava e você não sabe qual addon é o responsável,
\ui{Encontrar o culpado} procura por metades. A cada rodada, o programa ativa
parte dos seus addons, você abre o simulador e responde \ui{Travou} ou
\ui{Funcionou}, e metade dos suspeitos que restam é descartada.

\shot{31-diagnostics-bisection.png}{A busca pelo culpado, com a rodada atual, o
que já foi descartado e as duas respostas.}

Antes de começar:

\begin{itemize}
  \item A busca trabalha com unidades. Uma unidade é um addon, ou um grupo de
    addons que precisam ficar juntos, por exemplo porque dividem uma pasta de
    modelo. O painel diz por que um grupo está junto, e oferece dividi-lo quando
    isso é seguro.
  \item A primeira rodada desativa todos os seus addons, para saber se a causa
    está mesmo entre eles.
  \item A sua configuração é guardada antes da primeira rodada e restaurada
    quando a busca terminar, do jeito que terminar. Se o programa fechar durante
    a busca, ele oferece restaurar a configuração na próxima vez que abrir.
  \item A busca supõe um único culpado. Se a falha só acontece com dois addons
    ativados juntos, o resultado pode apontar o errado.
\end{itemize}

\section{Histórico}

Toda operação que altera arquivos no disco acrescenta uma linha ao
\ui{Histórico}. As linhas só são acrescentadas, e o programa nunca apaga o
histórico.

Cada linha mostra quando aconteceu, a operação, o addon, a biblioteca, a origem,
o destino e o resultado. \ui{Só falhas} mostra só o que deu errado, e a busca
procura em nomes de addon, caminhos e operações.

O \ui{Desfazer última alteração}, na aba \ui{Biblioteca}, desfaz a última
alteração inteira, e não uma linha dela.

\section{Opções}

O botão de engrenagem no cabeçalho abre as opções, que têm cinco abas.

\shot{10-options-profiles.png}{A aba de perfis e bibliotecas.}

\begin{description}
  \item[Perfis e bibliotecas] Escolha o perfil em uso, adicione ou remova perfis
    e bibliotecas, e troque destinos. Remover uma biblioteca não exclui nenhum
    arquivo: os links dela continuam funcionando no simulador, mas o FS
    Organizer deixa de gerenciá-los. \ui{Categorias} mostra quantas pastas de
    categoria estão marcadas e deixa marcar ou remover os marcadores.
  \item[Links] O tipo de link, e como as importações são conferidas.
  \item[Atualizações] Automática, avisar ou manual. A automática baixa as
    versões novas e as instala quando você fecha o programa.
  \item[Idioma] Inglês ou português do Brasil. A troca vale na hora.
  \item[Sobre] A versão, a licença e os créditos.
\end{description}

\section{Onde ficam as coisas}

O programa guarda os dados dele em \where|%LOCALAPPDATA%\fs-organizer|.
\where|settings.json| guarda os perfis, as bibliotecas, os destinos e as suas
escolhas. \where|journal\operations.jsonl| é o histórico. \where|presets|
guarda um arquivo por preset, mais o preset de retorno de cada perfil.
\where|bisection| guarda o estado de uma busca pelo culpado em andamento, e é
apagado quando a busca termina.

Nenhum addon seu fica ali. Se você apagar essa pasta, perde a configuração e os
presets, mas nenhum addon: o programa os encontra de novo lendo os destinos e as
bibliotecas.

A única coisa que o programa grava fora dessa pasta é o marcador oculto
\where|.fsorg-category| nas pastas de categoria. \ui{Categorias}, nas opções,
remove todos eles.

\end{document}
